\documentclass[11pt]{article}

\usepackage[margin=1in]{geometry}
\usepackage{amsmath,amssymb,amsthm}
\usepackage{bm}
\usepackage{hyperref}
\usepackage{parskip}

\renewcommand{\ne}{n_{e}}  % shadow LaTeX kernel's \ne ("not equal")
\newcommand{\nfc}{n_{\mathrm{FC}}}
\newcommand{\ncx}{n_{\mathrm{CX}}}
\newcommand{\Vfc}{\lvert V_{\mathrm{FC}}\rvert}
\newcommand{\Vcx}{\lvert V_{\mathrm{CX}}\rvert}
\newcommand{\Si}{S_{i}}
\newcommand{\Scx}{S_{\mathrm{CX}}}
\newcommand{\Dped}{D_{\mathrm{ped}}}
\newcommand{\fFC}{f_{\mathrm{FC}}}
\newcommand{\fCX}{f_{\mathrm{CX}}}
\newcommand{\fsa}[1]{\left\langle #1 \right\rangle}
\newcommand{\gradrsq}{\fsa{|\nabla r|^{2}}}
\newcommand{\neginf}{x=-\infty}
\newcommand{\dnein}{\left.\frac{d\ne}{dx}\right|_{\mathrm{in}}}
\newcommand{\xin}{x_{\mathrm{inner}}}

\title{Derivation of the \emph{complete} form of report Eq.~(6)\\
  \large (\texttt{eq6\_form="complete"} in \texttt{src/solver\_sc.py})}
\author{}
\date{}

\begin{document}
\maketitle

\begin{abstract}
\noindent
Equation~(6) of the report is the no-charge-exchange ``first step'' of the
Saarelma--Connor pedestal model --- Eq.~(16) of Saarelma \emph{et al.}~2023,
Nucl.\ Fusion \textbf{63} 052002.  This note derives the form that the code
actually solves by default, \texttt{eq6\_form="complete"}, states exactly
which two approximations separate it from the equation as printed in the
paper, and carries it through to the non-dimensional coefficients
implemented in \texttt{solve\_sc\_scipy} and \texttt{solve\_sc\_firedrake}.
The companion document \texttt{derivation\_eq16.tex} contains the full
chain from the three-fluid model; this note is the code-facing summary of
its Eq.~(16-full), plus the two pieces that document does not cover: the
$\ne$-dependent conductance and the numerical reason the complete form is
the default.
\end{abstract}

\section{Starting point}
\label{sec:start}

The derivation in \texttt{derivation\_eq16.tex} carries the three-fluid
model Eqs.~(1)--(3) through flux-surface averaging (Eqs.~4--6), the form
factors (Eq.~7), the two neutral kinetic balances (Eqs.~8--9) and one
radial integration (Eq.~10).  Two results from that chain are all we need
here.

\paragraph{(i) The electron transport equation, flux-surface averaged.}
\begin{equation}
  \frac{d}{dx}\!\left[\gradrsq\,\Dped\,\frac{d\ne}{dx}\right]
    \;=\; -\,\ne\,\Si\,\bigl(\fsa{\nfc}+\fsa{\ncx}\bigr),
  \label{eq:transport}
\end{equation}
with $x = r - r_{\mathrm{sep}} < 0$ inside the pedestal.

\paragraph{(ii) The once-integrated total-flux relation (Eq.~10).}
\begin{equation}
  \gradrsq\,\Dped\,\frac{d\ne}{dx}
    \;=\;
    -\,\frac{1+\Scx/(2\Si)}{1+\Scx/\Si}\,\Vfc\,\fFC\,\fsa{\nfc}
    \;-\; \Vcx\,\fCX\,\fsa{\ncx}
    \;+\; C,
  \label{eq:eq10}
\end{equation}
where the integration constant is fixed by matching to the core, where both
neutral densities vanish:
\begin{equation}
  C \;=\; \left[\gradrsq\,\Dped\,\frac{d\ne}{dx}\right]_{\neginf}
      \;\simeq\; \gradrsq\,\Dped\,\dnein .
  \label{eq:C}
\end{equation}
(The code imposes exactly this as the inner Neumann boundary condition,
$\ne'(\xin) = \dnein$, so $C\neq 0$ in general --- unlike the original
paper, which sets $C\equiv0$.)

The prefactor $\bigl[1+\Scx/(2\Si)\bigr]/\bigl[1+\Scx/\Si\bigr]$
in~\eqref{eq:eq10} is the \emph{shape factor}.  It arises from summing the
FC and CX source terms in~\eqref{eq:transport}: the FC balance contributes
$\Si/(\Si+\Scx)$ and the CX return flux contributes $\Scx/[2(\Si+\Scx)]$,
and
\begin{equation}
  \frac{\Si}{\Si+\Scx} + \frac{\Scx}{2(\Si+\Scx)}
  \;=\; \frac{2\Si+\Scx}{2(\Si+\Scx)}
  \;=\; \frac{1+\Scx/(2\Si)}{1+\Scx/\Si}.
  \label{eq:shape}
\end{equation}
The $\tfrac12$ is the half of the outgoing CX flux that is isotropised back
inward.  Everything below turns on whether this factor is kept.

\section{The no-CX closure}
\label{sec:closure}

\textbf{Assumption A6 (no CX neutrals):} set $\fsa{\ncx}\to0$.  This is the
whole content of the ``first step'' --- it makes the equation solvable in
one shot and gives the Picard loop for the full Eq.~(7) a physically
sensible starting profile.  Under A6, Eq.~\eqref{eq:eq10} becomes an
algebraic relation for $\fsa{\nfc}$:
\begin{equation}
  \fsa{\nfc}
    \;=\; \frac{1+\Scx/\Si}{1+\Scx/(2\Si)}\;
          \frac{C \;-\; \gradrsq\,\Dped\,\dfrac{d\ne}{dx}}{\Vfc\,\fFC}.
  \label{eq:closure}
\end{equation}
This is the closure: the FC neutral density is not an unknown any more, it
is slaved to the local electron flux.  Inserting~\eqref{eq:C} for $C$ and
using assumption~A2 ($\gradrsq$ slowly varying, so the same $\gradrsq\Dped$
multiplies both slopes),
\begin{equation}
  \fsa{\nfc}
    \;=\; \frac{1+\Scx/\Si}{1+\Scx/(2\Si)}\;
          \frac{\gradrsq\,\Dped}{\Vfc\,\fFC}\,
          \left(\dnein - \frac{d\ne}{dx}\right).
  \label{eq:closure2}
\end{equation}
Note the sign: inside the pedestal $d\ne/dx > \dnein$ (the profile steepens
towards the separatrix), so $\fsa{\nfc}>0$ as it must be.

\section{Substitution: the complete report Eq.~(6)}
\label{sec:complete}

Put~\eqref{eq:closure2} into the no-CX form of the transport
equation~\eqref{eq:transport}, $\frac{d}{dx}[\gradrsq\Dped\ne'] =
-\ne\Si\fsa{\nfc}$:
\begin{align}
  \frac{d}{dx}\!\left[\gradrsq\,\Dped\,\frac{d\ne}{dx}\right]
    &= -\,\ne\,\Si\,\frac{1+\Scx/\Si}{1+\Scx/(2\Si)}\,
         \frac{\gradrsq\,\Dped}{\Vfc\,\fFC}\,
         \left(\dnein - \frac{d\ne}{dx}\right).
\end{align}
Flip the sign of the last bracket and use
$\Si\,(1+\Scx/\Si) = \Si+\Scx$:
\begin{equation}
  \boxed{\;
  \frac{d}{dx}\!\left[\gradrsq\,\Dped\,\frac{d\ne}{dx}\right]
    \;=\; \frac{\ne\,(\Si+\Scx)}{1+\Scx/(2\Si)}\;
          \frac{\gradrsq\,\Dped}{\Vfc\,\fFC}\;
          \left(\frac{d\ne}{dx} - \dnein\right).
  \;}
  \label{eq:6complete}
\end{equation}
This is the complete form --- Eq.~(16-full) of
\texttt{derivation\_eq16.tex}, and the equation selected by
\texttt{eq6\_form="complete"}.  Every geometric factor and every form
factor is still explicit; no cancellation has been performed.

\subsection{What separates it from the equation as printed}
\label{sec:vs-paper}

The paper writes
\begin{equation}
  \frac{d}{dx}\!\left[\gradrsq\,\Dped\,\frac{d\ne}{dx}\right]
    \;=\; \ne\,\Dped\,\frac{\Si+\Scx}{\Vfc\,\fFC}\,
          \left(\frac{d\ne}{dx} - \dnein\right),
  \label{eq:6paper}
\end{equation}
which is \texttt{eq6\_form="paper"}.  Exactly two steps take
\eqref{eq:6complete} to \eqref{eq:6paper}:

\begin{itemize}
  \item \textbf{A5 --- drop the shape factor.}  Approximate
        $1+\Scx/(2\Si)\approx1$, i.e.\ $\Scx\ll2\Si$.  In the pedestal
        $\Scx$ is comparable to $\Si$ (often larger at low $T_e$), so this
        is the weakest link in the chain; the complete form simply keeps
        the factor, which costs nothing.
  \item \textbf{Symbolic cancellation of $\gradrsq$.}  The numerator
        $\gradrsq\,\Dped$ of the closure is read as ``$\Dped$ times a
        $\gradrsq$ that cancels against the $\gradrsq$ of the transport
        operator on the left''.  That cancellation is legitimate only if
        $\gradrsq\approx1$; it is a third, unstated $\gradrsq\to1$
        approximation \emph{inside} the neutral closure, not an algebraic
        identity.  In~\eqref{eq:6complete} the two $\gradrsq\Dped$ factors
        sit on opposite sides of the equation and cancel \emph{exactly}
        when the equation is divided through by the conductance
        (Sec.~\ref{sec:expanded}), whatever the value of $\gradrsq$.
\end{itemize}

The two forms therefore differ by the multiplicative factor
\begin{equation}
  \frac{\text{paper RHS}}{\text{complete RHS}}
    \;=\; \frac{1+\Scx/(2\Si)}{\gradrsq}.
  \label{eq:ratio}
\end{equation}

\paragraph{Why this is not a cosmetic difference.}
On the DIII-D 158091 case shipped with this repository
$\gradrsq \sim 0.05$, so~\eqref{eq:ratio} is ${\sim}35$.  The first-step
equation~\eqref{eq:6complete} is a
Bernoulli-type equation whose solution amplifies roughly like
$\exp\bigl[\int C_{A6}\,N\,d\xi\bigr]$ across the pedestal
(Sec.~\ref{sec:nondim}); a factor 35 in $C_{A6}$ turns an amplification of
$\sim\!70$ into $\exp(176)$.  The complete form converges in both
implementations; the printed form is not solvable in double precision on
this equilibrium at all, and the solver reports the failure through
\texttt{\_first\_step\_failed\_sc}.  Hence \texttt{"complete"} is the
default.

\paragraph{Consistency with Eq.~(7).}  The complete form is also the one
that matches the second-step equation, report Eq.~(7), which keeps
$\gradrsq\Dped$ intact in exactly the same slot:
\[
  \frac{d}{dx}\!\left[\gradrsq\Dped\frac{d\ne}{dx}\right]
   = -\ne\Si\left\{-\frac{\gradrsq\Dped}{\Vcx\,\fCX}
     \left(\frac{d\ne}{dx}-\dnein\right)
     + C_{\mathrm{CX}}(x)\,\fsa{\nfc}(0)\,E(x)\right\}.
\]
Using \texttt{"paper"} for step~1 and Eq.~(7) for step~2 means starting the
Picard loop from the solution of a differently-scaled equation.

\paragraph{On $\fFC$.}  The report prints Eq.~(6) with $\fFC=1$
(assumption A4), which is what \texttt{form\_factor} in
\texttt{saarelma\_connor} currently sets.  Both code paths carry $\fFC$ explicitly, so for the
present form-factor model the $\fFC$ placement is immaterial; it is kept
symbolic so a non-trivial poloidal model can be dropped in later.

\section{Expanded (non-conservative) form with an $\ne$-dependent conductance}
\label{sec:expanded}

The scipy path solves the equation in expanded form, which requires care
because the pedestal conductance depends on $\ne$ itself.  Write
\begin{equation}
  f(x,\ne) \;\equiv\; \gradrsq\,\Dped
    \;=\; \underbrace{\gradrsq\,(D_{\mathrm{NEO}}+D_{\mathrm{KBM}})}_{\textstyle f_{0}(x)}
      \;+\; \frac{\overbrace{\gradrsq\,C_{\mathrm{ETG}}}^{\textstyle f_{1}(x)}}{\ne},
  \label{eq:fsplit}
\end{equation}
the split implemented by \texttt{\_build\_f\_interp} /
\texttt{\_conductance}: the ETG channel enters as a heat-flux-like term
$\propto C_{\mathrm{ETG}}/\ne$, so $f$ is a function of both $x$ and $\ne$.
Expanding the divergence with the chain rule,
\begin{equation}
  \frac{d}{dx}\bigl[f\,\ne'\bigr]
    \;=\; f\,\ne'' \;+\; \frac{\partial f}{\partial x}\,\ne'
      \;+\; \frac{\partial f}{\partial \ne}\,(\ne')^{2},
  \qquad
  \frac{\partial f}{\partial x} = f_{0}' + \frac{f_{1}'}{\ne},
  \quad
  \frac{\partial f}{\partial \ne} = -\frac{f_{1}}{\ne^{2}}.
\end{equation}
Insert this on the left of~\eqref{eq:6complete} and divide by $f$.  The
$\gradrsq\Dped = f$ of the right-hand side cancels the $f$ of the division
\emph{exactly} --- this is the cancellation that the printed form performs
prematurely and only approximately:
\begin{equation}
  \boxed{\;
  \ne'' \;=\; c_{A6}(x)\,\ne\,\bigl(\ne' - \dnein\bigr)
          \;-\; c_{K}\,\ne' \;+\; c_{N}\,(\ne')^{2},
  \;}
  \label{eq:phys-ode}
\end{equation}
with
\begin{equation}
  c_{A6} = \frac{\Si+\Scx}{\bigl[1+\Scx/(2\Si)\bigr]\,\Vfc\,\fFC},
  \qquad
  c_{K} = \frac{1}{f}\left(f_{0}' + \frac{f_{1}'}{\ne}\right),
  \qquad
  c_{N} = \frac{f_{1}}{\ne^{2} f}.
\end{equation}
Note that $c_{A6}$ carries no $\gradrsq$ at all.  For the printed form the
same reduction leaves
$c_{A6}^{\text{paper}} = (\Si+\Scx)/(\gradrsq\,\Vfc\,\fFC)$ --- the
$1/\gradrsq$ of Eq.~\eqref{eq:ratio} survives because there the right-hand
side carried only $\Dped$, not $\gradrsq\Dped$, so the division by $f$ left
a residual $1/\gradrsq$.

\section{Non-dimensionalization}
\label{sec:nondim}

Use the scalings of \texttt{solver\_sc.py},
\begin{equation}
  \xi = \frac{x}{L}\in[-1,0],\quad L=|\xin|,
  \qquad
  N = \frac{\ne}{n_{0}},\quad n_{0}=\texttt{ne\_x0},
  \qquad
  N'_{\mathrm{in}} = \frac{L}{n_{0}}\dnein,
\end{equation}
so that $\ne' = (n_{0}/L)N'$ and $\ne'' = (n_{0}/L^{2})N''$.  Substituting
into~\eqref{eq:phys-ode} and multiplying through by $L^{2}/n_{0}$:
\begin{alignat}{2}
  \ne'' &\;\to\; \frac{n_{0}}{L^{2}}N''
    &&\;\Longrightarrow\; N'', \\
  c_{A6}\ne(\ne'-\dnein) &\;\to\; c_{A6}\,n_{0}N\,\frac{n_{0}}{L}(N'-N'_{\mathrm{in}})
    &&\;\Longrightarrow\; \bigl(n_{0}L\,c_{A6}\bigr)N(N'-N'_{\mathrm{in}}), \\
  c_{K}\ne' &\;\to\; c_{K}\frac{n_{0}}{L}N'
    &&\;\Longrightarrow\; \bigl(L\,c_{K}\bigr)N', \\
  c_{N}(\ne')^{2} &\;\to\; c_{N}\frac{n_{0}^{2}}{L^{2}}(N')^{2}
    &&\;\Longrightarrow\; \bigl(n_{0}c_{N}\bigr)(N')^{2}.
\end{alignat}
Hence
\begin{equation}
  \boxed{\;
    N'' \;=\; C_{A6}\,N\,(N' - N'_{\mathrm{in}}) \;-\; C_{K}\,N' \;+\; C_{N}\,(N')^{2},
  \;}
\end{equation}
\begin{equation}
  \boxed{\;
  \begin{aligned}
    C_{A6} &\;=\; \frac{n_{0}L\,(\Si+\Scx)}{\bigl[1+\Scx/(2\Si)\bigr]\,\Vfc\,\fFC}
      \qquad\text{(\texttt{"complete"})}, \\[2pt]
    C_{A6} &\;=\; \frac{n_{0}L\,(\Si+\Scx)}{\gradrsq\,\Vfc\,\fFC}
      \qquad\qquad\ \ \text{(\texttt{"paper"})}, \\[4pt]
    C_{K} &\;=\; \frac{L}{f}\left(\frac{df_{0}}{dx} + \frac{1}{n_{0}N}\frac{df_{1}}{dx}\right),
    \qquad
    C_{N} \;=\; \frac{f_{1}}{n_{0}N^{2}f}.
  \end{aligned}
  \;}
\end{equation}
These are the coefficients assembled in \texttt{ode\_first} inside
\texttt{solve\_sc\_scipy}, with $C_{K}$ and $C_{N}$ returned by
\texttt{\_conductance} and $f_{0},f_{1}$ and their derivatives supplied by
the shape-preserving PCHIP interpolants of \texttt{\_build\_f\_interp}
($f$ needs a derivative, so a PCHIP is used rather than the linear
\texttt{np.interp} used elsewhere in the repository).

\paragraph{Amplification estimate.}  Neglecting $C_{K}$ and $C_{N}$ and
writing $u = N' - N'_{\mathrm{in}}$, Eq.~\eqref{eq:phys-ode} reads
$u' \simeq C_{A6}N\,u$, so $u$ grows like
$\exp\bigl[\int_{-1}^{\xi}C_{A6}N\,d\xi'\bigr]$.  This is the estimate quoted in
Sec.~\ref{sec:vs-paper}: the exponent is directly proportional to
$C_{A6}$, hence to $\bigl[1+\Scx/(2\Si)\bigr]/\gradrsq$ when comparing the
two forms.

\section{Conservative (weak) form}
\label{sec:weak}

The Firedrake path keeps~\eqref{eq:6complete} in conservative form.  With
$[V]_{0} = Ln_{0}S_{0}$, $[D]_{0} = L^{2}n_{0}S_{0}$, $S_{0}=\Si(0)$, and
hatted quantities scaled accordingly ($\hat f = f/[D]_{0}$,
$\hat S = S/S_{0}$, $\hat V = V/[V]_{0}$), the non-dimensional strong form
is
\begin{equation}
  \frac{d}{d\xi}\!\left[\hat f\,N'\right]
    \;=\; N\,D_{6}\,\bigl(\hat S_{i}+\hat S_{\mathrm{CX}}\bigr)\,
          \frac{N' - N'_{\mathrm{in}}}{\hat V_{\mathrm{FC}}\,\fFC},
  \qquad
  D_{6} \;=\;
  \begin{cases}
    \dfrac{\hat f}{1+\hat S_{\mathrm{CX}}/(2\hat S_{i})}, & \text{\texttt{"complete"}},\\[10pt]
    \dfrac{\hat f}{\gradrsq}, & \text{\texttt{"paper"}},
  \end{cases}
  \label{eq:strong-nd}
\end{equation}
which is precisely the \texttt{D6} switch in
\texttt{solve\_sc\_firedrake}.  The two branches again differ by the factor
$\bigl[1+\hat S_{\mathrm{CX}}/(2\hat S_{i})\bigr]/\gradrsq$ of
Eq.~\eqref{eq:ratio}; note that this ratio is scale-invariant, so it is the
same in hatted and unhatted variables.

Multiplying~\eqref{eq:strong-nd} by a test function $w$ and integrating by
parts over $\xi\in[-1,0]$,
\begin{equation}
  \int_{-1}^{0}\hat f\,N'\,w'\,d\xi
  \;+\; \int_{-1}^{0} N\,D_{6}\,(\hat S_{i}+\hat S_{\mathrm{CX}})
        \frac{N'-N'_{\mathrm{in}}}{\hat V_{\mathrm{FC}}\fFC}\,w\,d\xi
  \;+\; \bigl[\hat f\,N'_{\mathrm{in}}\,w\bigr]_{\xi=-1}
  \;=\; 0,
\end{equation}
with $w$ vanishing at $\xi=0$ where the Dirichlet condition $N(0)=1$ is
imposed strongly.  The boundary term at $\xi=-1$ \emph{is} the inner
Neumann condition $N'(-1)=N'_{\mathrm{in}}$, imposed naturally as a
prescribed flux (\texttt{bnd\_term} on boundary id 1); the $\hat f$
multiplying it is the same $\ne$-dependent conductance
of~\eqref{eq:fsplit}, assembled inline as
$\hat f = g(\hat D_{\mathrm{NEO}}+\hat D_{\mathrm{KBM}}) + g\hat
C_{\mathrm{ETG}}/N$ so that Newton sees the $1/N$ nonlinearity.
Because this path never divides by $\hat f$, the cancellation discussed in
Sec.~\ref{sec:expanded} is not performed symbolically at all --- it happens
in the assembled residual, which is why the two implementations agree to
$\sim\!3\times10^{-3}$ in max-norm on the converged profile.

\section{Boundary conditions and consistency}
\label{sec:bcs}

Equation~\eqref{eq:6complete} is second order and needs two conditions.
The only supported combination is
\begin{equation}
  \left.\frac{d\ne}{dx}\right|_{x=\xin} = \dnein
  \quad\text{(inner, Neumann)},
  \qquad
  \ne(0) = \texttt{ne\_x0}
  \quad\text{(separatrix, Dirichlet)}.
\end{equation}
The Neumann condition is not an independent modelling choice: it is the
same slope that fixed the integration constant $C$ in~\eqref{eq:C}, so
using any other value at $\xin$ would make the closure~\eqref{eq:closure2}
inconsistent with the equation being solved.  Note the built-in
consistency check: at $x=\xin$ the right-hand side
of~\eqref{eq:6complete} vanishes identically, so the flux
$\gradrsq\Dped\ne'$ is stationary there --- as it must be, since
$\fsa{\nfc}(\xin)=0$ means no neutral source at the inner boundary.

A converged solution with $\ne\le0$ anywhere is rejected by both paths.
Equation~\eqref{eq:6complete} admits such solutions (nothing in it enforces
positivity), but they are unphysical and no Eq.~(7) Picard iterate can
recover from them.

\section{Dimensional check}
\label{sec:dimcheck}

With $[\Dped]=[f]=\mathrm{m^{2}\,s^{-1}}$, $[S]=\mathrm{m^{3}\,s^{-1}}$,
$[\Vfc]=\mathrm{m\,s^{-1}}$, $[n_{0}]=\mathrm{m^{-3}}$, $[L]=\mathrm{m}$,
and $\gradrsq$, $\fFC$, $\Scx/\Si$ all dimensionless:
\begin{equation}
  [C_{A6}] = \frac{\mathrm{m^{-3}}\cdot\mathrm{m}\cdot\mathrm{m^{3}\,s^{-1}}}
                  {\mathrm{m\,s^{-1}}} = 1,
  \qquad
  [C_{K}] = \frac{\mathrm{m}\cdot\mathrm{m\,s^{-1}}}{\mathrm{m^{2}\,s^{-1}}} = 1,
  \qquad
  [C_{N}] = \frac{\mathrm{m^{-1}\,s^{-1}}}
                 {\mathrm{m^{-3}}\cdot\mathrm{m^{2}\,s^{-1}}} = 1 .
\end{equation}
For $C_{N}$ note $[f_{1}] = [\gradrsq C_{\mathrm{ETG}}] =
\mathrm{m^{-1}\,s^{-1}}$ since $f_{1}/\ne$ must have units of $f$.  All
three coefficients are dimensionless.\hfill$\checkmark$

Both branches of $C_{A6}$ are dimensionless, so the dimensional check does
\emph{not} discriminate between them --- the difference is a pure number,
which is exactly why the discrepancy is easy to miss when reading the
paper.

\section{Assumptions used here}
\label{sec:assumptions}

Referenced by the labels of \texttt{derivation\_eq16.tex}:

\begin{itemize}
  \item[\textbf{A1}] Locally mono-energetic neutrals;
        $\Vfc=\sqrt{8E_{\mathrm{FC}}/(\pi^{2}M_{i})}$ constant.
        (Only the FC species appears in this note.)
  \item[\textbf{A2}] Slowly varying geometry: $\gradrsq$, $\oint R^{2}d\theta$
        and $\fFC$ pass through $d/dx$.  Used twice: in deriving
        Eq.~\eqref{eq:eq10}, and in writing~\eqref{eq:closure2} with the
        same $\gradrsq\Dped$ on both slopes.
  \item[\textbf{A3}] Slab-like neutral closure: one explicit $\gradrsq$ is
        dropped in the FC kinetic balance, while the transport
        operator~\eqref{eq:transport} retains its own.
  \item[\textbf{A4}] $\fFC=1$ (poloidally symmetric FC neutrals) --- what
        \texttt{form\_factor} currently returns.  Carried symbolically
        in~\eqref{eq:6complete} and in both implementations.
  \item[\textbf{A6}] No CX neutrals, $\fsa{\ncx}=0$.  Defines the first
        step; lifted iteratively by the Eq.~(7) Picard loop.
  \item[\textbf{A7}] Inner boundary $\ne'(\xin)=\dnein$ with
        $\fsa{\nfc}(\xin)\approx0$, giving the finite $C$ of~\eqref{eq:C}.
  \item[\textbf{A8}] Quasineutrality with $Z_{i}=1$, $n_{i}\approx\ne$.
\end{itemize}

\textbf{Not} used --- this is the point of the complete form:
\begin{itemize}
  \item[\textbf{A5}] $1+\Scx/(2\Si)\approx1$ (dropping the shape factor);
  \item[] the unstated $\gradrsq\approx1$ inside the neutral closure that
        licenses cancelling $\gradrsq\Dped$ against $\Dped$.
\end{itemize}
Restoring both recovers the paper's Eq.~(16), i.e.\
\texttt{eq6\_form="paper"}.

\end{document}
